\addcontentsline{toc}{section}{Podsumowanie i wnioski końcowe}

Niniejsza praca magisterska stanowiła kompleksowe studium problemu uczenia (aproksymacji) funkcji wysokowymiarowych zależących od niewielkiej liczby ukrytych parametrów liniowych. Głównym wyzwaniem badawczym było ominięcie fundamentalnej bariery, jaką w klasycznej analizie numerycznej stanowi klątwa wymiarowości \cite{novak2008tractability}. 

Opierając się na uogólnionym modelu funkcji $f(x) = g(Ax)$ wprowadzonym w pracy \cite{fornasier2012learning}, wykazaliśmy, że dzięki wykorzystaniu struktury rzadkości gradientów kierunkowych możliwe jest skonstruowanie efektywnych algorytmów uczenia. Kombinacja aparatu teorii próbkowania skompresowanego (ang. \textit{compressed sensing}), zaawansowanych macierzowych nierówności probabilistycznych \cite{tropp2012user} oraz twierdzeń o stabilności rozkładu SVD \cite{wedin1972perturbation} pozwoliła na rygorystyczne udowodnienie, że błąd aproksymacji może być ograniczony z niezwykle wysokim prawdopodobieństwem przy użyciu liczby próbek rosnącej co najwyżej wielomianowo z wymiarem przestrzeni.

Przeprowadzone w ramach pracy doświadczenia numeryczne jednoznacznie potwierdziły poprawność teoretycznych oszacowań. Symulacje dla wielowymiarowych funkcji testowych wykazały, że:
\begin{enumerate}
    \item \textbf{Przełamanie klątwy wymiarowości:} Złożoność próbkowania, niezbędna do stabilnej rekonstrukcji wektorów i macierzy parametrów, rośnie logarytmicznie ($m_\Phi \approx \mathcal{O}(\log d)$) względem wymiaru przestrzeni $d$, co czyni metodę traktowalną nawet dla $d \ge 1000$.
    \item \textbf{Stabilność ekstrakcji podprzestrzeni:} Algorytm 2, wykorzystujący dekompozycję SVD macierzy różnic skończonych, skutecznie odnajduje podprzestrzeń aktywnych kierunków dzięki wyraźnej i stabilnej przerwie spektralnej.
    \item \textbf{Wrażliwość na parametry i szum:} Istnieje istotny kompromis przy doborze kroku różnicy skończonej $\varepsilon$. Podczas gdy w środowisku idealnym mniejsze $\varepsilon$ gwarantuje wyższą precyzję dzięki lepszemu przybliżeniu ze wzoru Taylora, w obecności szumu pomiarowego prowadzi ono do drastycznej amplifikacji błędu (tzw. \textit{noise folding}). Zwiększenie tolerancji na szum wymaga zastosowania większego $\varepsilon$ oraz relaksacji optymalizacyjnej (np. BPDN lub estymatora Dantzig selector \cite{candes2007dantzig}).
\end{enumerate}

Pomimo udowodnionej skuteczności, zaprezentowane metody otwierają szereg pytań badawczych na przyszłość. Istotnym ograniczeniem pozostaje skalowanie metody dla dowolnych ciał wypukłych (np. hiperkostek $[-1, 1]^d$), gdzie pojawiający się w oszacowaniach błędu czynnik średnicy dziedziny ponownie wprowadza bezpośrednią zależność od wymiaru $d$. Pokonanie tej przeszkody może w przyszłości wymagać zastosowania estymat opartych na funkcjonałach Minkowskiego i normach dualnych. Ponadto, wysoce obiecującym kierunkiem rozwoju uogólnionego modelu ($k \ge 1$) wydaje się bezpośrednie poszukiwanie macierzy rzadkiej poprzez minimalizację normy nuklearnej \cite{fornasier2010numerical}, zamiast kolumnowej minimalizacji normy $l_1$.

Podsumowując, zintegrowanie metod geometrii wysokich wymiarów z algorytmami optymalizacji wypukłej dostarcza niezwykle potężnych i stabilnych numerycznie narzędzi, które z powodzeniem mogą stanowić fundament dla nowych architektur uczenia maszynowego i kompresji danych, operujących w środowiskach o ogromnej liczbie cech.